\section{External results and the order of choices}
\label{app:dependencies}

The proof is self-contained except for the results listed in
Table~\ref{tab:external}. The table records each of them in the form
in which it is used, so that the reader can see exactly what is taken
from the literature. The rigidity and flexibility theorems quoted in
the introduction give context only; none of their conclusions is used.

\renewcommand{\arraystretch}{1.12}\setlength{\tabcolsep}{4pt}
\begin{longtable}{@{}p{.21\linewidth}p{.47\linewidth}p{.25\linewidth}@{}}
\caption{External results in the form used.}\label{tab:external}\\
\toprule
Result & Statement used & Where\\
\midrule
\endfirsthead
\toprule
Result & Statement used & Where\\
\midrule
\endhead
Gauss equation and structure equations &
 For a surface with metric $g$ and second fundamental form
 $\mathrm{II}$ in a three-manifold, $K_g=\overline K+\det_g\mathrm{II}$,
 with $\overline K$ the sectional curvature of the tangent plane. For
 an oriented orthonormal coframe with
 $d\eta_1=\alpha\wedge\eta_2$, $d\eta_2=-\alpha\wedge\eta_1$:
 $d\alpha=-K\,\eta_1\wedge\eta_2$. &
 Lemma~\ref{lem:darboux-invariant},
 Proposition~\ref{prop:scalar-core}\\
Primitives and coverings &
 A closed one-form on a simply connected domain is exact. A proper
 local diffeomorphism of the plane which is the identity outside a
 compact set is a diffeomorphism. &
 Lemma~\ref{lem:disk-development}\\
Ordinary differential equations &
 Existence, uniqueness and smooth dependence on data and parameters;
 Gronwall's inequality. &
 radial models, flows, Proposition~\ref{prop:obs-monotone-tail}\\
Quadratic forms \cite[Theorems~2.12, 2.13]{Teschl2009} &
 A closed quadratic form which is bounded below is the form of a
 unique self-adjoint operator. &
 the models of Section~\ref{sec:radial-models};
 Lemma~\ref{lem:radial-deep-inverse}\\
Spectral theorem \cite[Chapter~3]{Teschl2009} &
 For a self-adjoint operator $A$ and $z$ outside its spectrum,
 $\|(A-z)^{-1}\|=\dist(z,\sigma(A))^{-1}$. &
 Lemmas~\ref{lem:central-gap}, \ref{lem:morse-arithmetic}\\
Weyl's theorem \cite[Theorem~6.19]{Teschl2009} &
 If the difference of the resolvents of two self-adjoint operators
 is compact, their essential spectra coincide. &
 Section~\ref{subsec:morse}\\
Harmonic oscillator \cite[Section~8.3]{Teschl2009} &
 The operator $-\partial_x^2+\omega^2x^2$ on the line has the simple
 eigenvalues $(2n+1)\omega$, $n\ge0$, and no other spectrum. &
 Section~\ref{sec:worked-example},
 Lemma~\ref{lem:central-gap}\\
Countable products &
 A sequence in a countable product of compact intervals has a
 subsequence which converges in every coordinate. &
 Lemma~\ref{lem:radial-product-selection}\\
Three-lines lemma &
 A strongly analytic family of operators on $0<\operatorname{Re}z<1$,
 strongly continuous on the closed strip and bounded there on each
 vector, with operator norm bounds $M_0,M_1\ge0$ on its two boundary
 lines, is bounded by
 $M_0^{1-\vartheta}M_1^{\vartheta}$ on the line
 $\operatorname{Re}z=\vartheta$. &
 Lemma~\ref{lem:ang-base-inverse-regularity}\\
Extension operator &
 On a half-line or an interval there is a linear extension operator
 to the line which is bounded in every Sobolev norm. &
 after Lemma~\ref{lem:one-high-factor}\\
Hamilton's theorem on families of inverses &
 Theorem~\ref{thm:hamilton-inverse}. &
 Propositions~\ref{prop:rewrite-volume-inverse},
 \ref{prop:approx-inverse}\\
Hamilton's theorem with quadratic error &
 Theorem~\ref{thm:hamilton-quadratic}. &
 Theorem~\ref{thm:weak-branch}\\
\bottomrule
\end{longtable}

The Airy asymptotics \cite{DLMFAiry} quoted in
Section~\ref{sec:worked-example} serve as a heuristic only.

\paragraph{Two elementary facts used repeatedly.}
\emph{Schur's test.} If the blocks $K_{nm}$ of an operator matrix
satisfy $\sup_n\sum_m\|K_{nm}\|\le A$ and
$\sup_m\sum_n\|K_{nm}\|\le B$, then the operator is bounded on the
$\ell^2$ sum with norm at most $\sqrt{AB}$: apply the Cauchy--Schwarz
inequality with the weights $\|K_{nm}\|$ in each row, and then the
bound for the columns. \emph{Neumann series.} If $\|R\|<1$ on a
Banach space, then $\sum_{k\ge0}(-R)^k$ converges in norm to a
two-sided inverse of $I+R$. Whenever a Neumann series is used in this
paper, the space is one fixed Banach space and the bound $\|R\|<1$ is
established before any derivative of higher order is taken.

\paragraph{The order of the choices.}
The constants of the construction are chosen in the following order;
no later choice changes an earlier one.
\begin{enumerate}[label=\textup{\arabic*.}]
\item \emph{The models.} The parameters of the central oscillator and
 of the Morse models, the exponent $12$ of the joins and the
 constant $c_0$ of the matching angle
 (Sections~\ref{sec:radial-models}, \ref{sec:radial-tails}).
\item \emph{The profile.} The exponent $\ea$ of the source weight,
 then $C_*$, then the initial index $N$, then the cuts of the central
 model and the lengths of the cone heads; then the height ratios
 $\sigma_n$, by Lemma~\ref{lem:radial-product-selection}; then the
 connectors. These choices use only $|\gamma|\le1$ and are therefore
 the same for all endpoint exponents
 (Lemma~\ref{lem:radial-uniform-exponents}).
\item \emph{The reference inverse.} A compact interval of endpoint
 exponents $\gamma\in(0,1)$ changes the constants of the estimates,
 not the profile. The final border $(g_*,\ell_*)$ is fixed after the
 geometry (Section~\ref{subsec:final-border}).
\item \emph{The spaces.} The exponent $\beta=\tfrac34$, the collars
 $\delta_n^{\rm root}$, the controls $\chi_n$ and the cutoffs
 (Section~\ref{subsec:root-geometry}).
\item \emph{The neighbourhood.} The finite order $b_0$, then
 $\epsilon_0$, so that the comparisons of
 Lemma~\ref{lem:common-geometric-neighborhood} hold and the relative
 operator has norm less than $\tfrac12$ on one space
 (Lemma~\ref{lem:relative}).
\item \emph{The parameter.} The range of $q$, so that the control
 equation of Proposition~\ref{prop:approx-inverse} is solvable.
\end{enumerate}
Estimates of higher order change constants and use a fixed shift of
the order; they never require a new choice of the profile, of the
neighbourhood or of the range of $q$. The two inequalities
$\tfrac12<\beta<1$ serve different purposes: $\beta<1$ leaves room for
two neighbouring exponents at which the reference inverse is
available, and $\beta>\tfrac12$ makes the product of two decaying
factors decay faster than $H$.
